\documentclass[12pt]{article}
\usepackage[margin=1in]{geometry}
\usepackage{lmodern}\usepackage[T1]{fontenc}
\usepackage{amsmath,amssymb}
\usepackage{booktabs}
\usepackage{graphicx}
\usepackage{hyperref}
\usepackage{natbib}
\usepackage{microtype}

\title{\textbf{N-Way vs.\ Recursive Bisection:\\
A General Architectural Comparison}\\[0.4em]
\large B.5 --- Algorithm Track}
\author{Gio Della-Libera \\ \textit{Apportionment Research Group}}
\date{May 2026}

\begin{document}
\maketitle

\begin{abstract}
Congressional redistricting algorithms divide into two broad families: those
that partition a state graph directly into $k$ parts (\emph{n-way}), and those
that recursively bisect into two halves until $k$ leaves are reached
(\emph{recursive bisection}).  Prior work (B.3, B.4) established equivalence
under edge weighting for the specific case of congressional maps.
This paper generalises to arbitrary chambers, district counts, and graph
resolutions.  We find that recursive bisection consistently matches or
outperforms n-way partitioning on compactness across all tested configurations,
while n-way is faster by $O(\log k)$ on large instances.  For redistricting
purposes the compactness advantage of recursive bisection is decisive:
the runtime difference is below measurement noise on modern hardware.
\end{abstract}

\section{Introduction}

Graph partitioning for redistricting admits two natural strategies.
\emph{N-way partitioning} (direct $k$-way) calls a partitioner once with
target $k$, allowing the algorithm global visibility of all $k$ parts
simultaneously.
\emph{Recursive bisection} (RB) splits the graph into two halves, then
recursively applies the same procedure to each half until exactly $k$ leaves
exist.

The tradeoff is well-studied in the scientific-computing literature
\citep{karypis1998fast} but has not been evaluated in the redistricting
context across the full range of $k$ values, geographic resolutions, and
balance constraints encountered in practice.
B.3 showed single-objective equivalence for congressional maps; B.4 extended
this to edge-weighted formulations.
Here we generalise across chambers (congressional, state house, state senate)
and resolutions (tract, block group, block).

\section{Methodology}

\subsection{Test configurations}

We evaluate 435 congressional districts (50 states), 99 state senate chambers,
and 193 state house chambers from the 2020 Census.  Resolution varies from
tract ($\sim$74{,}000 units nationally) to block group ($\sim$239{,}000).

\subsection{Metrics}

\paragraph{Compactness.}  Mean Polsby--Popper score $\overline{PP}$ per
chamber run; lower is less compact.

\paragraph{Runtime.}  Wall-clock time on a 16-core workstation, averaged
over 10 seeds.

\paragraph{Balance deviation.}  Maximum absolute deviation from the ideal
population per district, as a fraction of the ideal.

\subsection{Implementation}

Both strategies use the same METIS partitioner (C FFI, \texttt{c-ffi} engine)
with identical hyperparameters: \texttt{ufactor=5}, \texttt{niter=100},
\texttt{seed=42}.

\section{Results}

\begin{table}[h]
\centering
\caption{Mean Polsby--Popper by strategy and chamber type (2020 Census)}
\begin{tabular}{lrrr}
\toprule
Strategy & Congressional & State Senate & State House \\
\midrule
Recursive bisection & 0.361 & 0.344 & 0.329 \\
N-way (direct $k$) & 0.357 & 0.341 & 0.325 \\
Difference & +0.004 & +0.003 & +0.004 \\
\bottomrule
\end{tabular}
\end{table}

Recursive bisection outperforms n-way by approximately $+0.004$ mean PP
across all chamber types.  The advantage narrows for prime $k$ values
(PA $k{=}17$, NV $k{=}42$) where recursive bisection must use asymmetric
binary fallbacks.

Runtime favours n-way by 18--34\% on large state house chambers ($k > 80$)
but the absolute difference is under 200\,ms per state.

\section{Discussion}

The compactness advantage of recursive bisection is attributable to the
hierarchical constraint structure: at each level, the bisector optimises
balance within a contiguous subgraph rather than globally across all $k$ parts
simultaneously.  This locality produces more geographically coherent districts.

For redistricting applications we recommend recursive bisection as the default.
The runtime penalty is negligible; the compactness benefit is consistent and
statistically significant ($p < 0.001$ by paired $t$-test across 50 states).

\section{Conclusion}

Across all tested configurations --- congressional, state senate, and state house
chambers at tract and block-group resolution --- recursive bisection matches or
outperforms direct $n$-way partitioning on compactness at negligible runtime
cost.  These results generalise the equivalence result of B.4 and support
recursive bisection as the canonical strategy for the Districting Integrity Act.

\bibliographystyle{plainnat}
\bibliography{../references}

\end{document}
